% Pursuing Stacks (cotes 134-2 to 134-8): archive cote 134-8, site page 8 (PDF page 9), typescript page 562.
% typescript. Draft compared: ps7.tex:350-423.
% AI-transcribed by Claude Opus 5.5 (claude-opus-5-5, adaptive thinking, effort high), Message Batches API, batch ps4 (msgbatch_012Aoc1NuskATPsyBeKSYK7J), 2026-11-12; reviewed against both supplied scan halves; see CHECK_DECISIONS.tsv.
\scanpage{134-8/8}
\mspage{562}

we use
\[ p_A = p_{A'} \hand{\circ} f \quad , \quad \xout{\text{hence }Lp} \]
hence
\[ \xout{\illegible{3}} \]
\[ \hand{(*)} \qquad (p_A)_!^{ab} \hand{\simeq} (p_{A'})_!^{ab} \hand{\circ} f_!^{ab} \quad , \quad \text{and} \quad L(p_A)_!^{ab} \hand{\simeq} L(p_{A'})_!^{ab} \hand{\circ} Lf_!^{ab} \quad , \]
taking into account that $f_!^{ab}$ maps projectives to projectives . Hence
we get
\[ LH.(A,F) \overset{def}{=} L(p_A)_!^{ab}(F) \hand{\simeq} L(p^{\deleted{\illegible{1}}}_{A'})_!^{ab}(Lf_!^{ab}(F)) \overset{def}{=} LH.(A',Lf_!^{ab}(F)) \unsure{.} \]
Thus, in order to get (5), we need only define a map in $D^-(A'{}^{\hand{\wedge}}_{ab})$
\[ (6) \qquad Lf_!^{ab}(F) \hand{\longrightarrow} F' \ , \]
which will be obtained as the composition
\[ (6') \qquad Lf_!^{ab}(F) \hand{\longrightarrow} f_!^{ab}(F) \hand{\longrightarrow} F' \ , \]
where the first map in (6') is the canonical augmentation map towards
the $H_o$ object, and where the second corresponds to u in (4) by adjuncti%
on. \over[typed]{T}{T}his defines the map (5), and compatibility with compositions
should be a tautology. Hence the functor (2). To get (1), we still
have to check that when
\[ f: \ A \hand{\longrightarrow} A' \]
is a weak equivalence, then
\[ LH.(f): \ LH.(A) \hand{\longrightarrow} LH.(A') \]
is an isomorphism in D.(Ab). It amounts to the same to check that for
any object $K^\bullet$ in $D^+(Ab)$, the corresponding map between the RHom's
with values in $K^\bullet$ is \xout{an quasi iso} an isomorphism in $D^+(Ab)$. B\over[typed]{u}{u}t \xout{the}
the latter map can be identified with the map for \uline{cohomology}
\[ RH^\bullet(A',K^\bullet_{A'} \ ) \hand{\longrightarrow} RH^\bullet(A, K^\bullet_A \ ) \quad , \]
with coefficients in the constant complex of presheaves defined by $K^\bullet$
on $A'$ and on $A$ , which is an isomorphism, by the very definition
of weak equivalences in (Cat) via cohomology.

Now the statement ``an integrator \xout{for A} is an abelianizator'' may
be rephrased rather evidently, without any reference to a given inte-%
grator, as \xout{bein} merely the commutativity, up to canonical isomorphism,
of the \xout{corres} following diagram for a given A in (Cat) :

\margin{8}

% Visual collation: both supplied image halves inspected; see CHECK_DECISIONS_134_7_8.tsv.
